\chapter{Faddeev-Popov Ghosts and BRST Symmetry}

\chapterquote{A field-theory formula is useful only after the smaller calculation inside it is visible.}

This chapter is part of \emph{Non-Abelian Theory, Broken Symmetry, and Advanced Topics}.  The working vocabulary is ghosts, gauge volume, BRST, physical states.  The first pass through the chapter should be slow: identify the object being changed, the condition held fixed, and the reason a boundary term or normalization is allowed.  The first concrete theme is deriving the non-Abelian gauge transformation; later sections reuse the same habit in a more compact notation.

For Faddeev-Popov Ghosts and BRST Symmetry, the calculus-level starting point is tied to deriving the non-Abelian gauge transformation.  Matrices, waves, operators, fields, and gauge variables are introduced only as the calculation requires them.  A formula is accepted after it has been checked in a finite model, differentiated or varied explicitly, and connected to the field-theoretic use that motivates it.

\section{The concrete problem: deriving the non-Abelian gauge transformation}

The work of this section is deriving the non-Abelian gauge transformation.  In the chapter heading the vocabulary is ghosts, gauge volume, BRST, physical states, but the first objects on the page are more prosaic: matrix-valued gauge fields, covariant field strengths, structure constants, ghosts, BRST symmetry, and self-interactions.  The formal notation will be useful only after it has been tied to a limit, a symmetry, a normalization condition, or an integration-by-parts step.  In Faddeev-Popov Ghosts and BRST Symmetry, section 1, the useful habit is to name the object, the operation, and the assumption before using the compact notation.

The finite model is a field with several components transformed by a position-dependent unitary matrix.  In that model no mystery is available: every derivative is an ordinary derivative with respect to one listed variable, every inner product is a sum, and every boundary assumption can be written at the endpoints.  This is why the finite model is not a toy in the dismissive sense.  It is the bookkeeping device that prevents continuum notation from becoming a fog of indices.  For Faddeev-Popov Ghosts and BRST Symmetry, section 1, that bookkeeping is deliberately modest, but it is what later keeps signs, factors of \(2\pi\), and normalization constants from turning into guesses.

The central idea for this chapter is the commutator of covariant derivatives produces a field strength that contains a gauge-field commutator.  To test that idea, strip the formula down until only the operation remains.  A sign change after raising or lowering an index means the metric has entered.  A derivative moving from one factor to another means a boundary term has been handled.  A normalization constant means that some unit object, such as a delta function, basis vector, or one-particle state, has been fixed.  Read in the setting of Faddeev-Popov Ghosts and BRST Symmetry, section 1, the line is a check on the calculation rather than an invitation to memorize another rule.

For the concrete problem: deriving the non-abelian gauge transformation, the calculation should also pass a dimensional check.  The left and right sides must carry the same units; more subtly, they must transform in the same way under the symmetries already introduced.  This is not a proof, but it is a quick way to catch wrong factors of \(2\pi\), signs in the metric, missing powers of the lattice spacing, or an omitted charge.  The same step will look more abstract later; in Faddeev-Popov Ghosts and BRST Symmetry, section 1 it is kept close to the finite calculation so the limiting move remains visible.

The bridge to quantum field theory is direct.  Yang-Mills theory supplies the strong and weak gauge sectors and introduces gauge-boson self-interactions.  Later notation will carry more indices and fewer verbal reminders, so the safe habit is to keep asking which operation is being performed and which quantity is being held fixed.  At this point in Faddeev-Popov Ghosts and BRST Symmetry, section 1, it is worth asking what would fail if the boundary condition, normalization, or symmetry assumption were changed.

One warning is worth making explicit here.  Ghost fields are not optional decoration in covariant gauges; they represent the Faddeev-Popov determinant.  This warning points to the delicate step of the derivation, not to a decorative aside.  In Faddeev-Popov Ghosts and BRST Symmetry, section 1, the calculation is meant to leave a trail from an ordinary derivative, sum, matrix product, or Gaussian integral to the field-theory formula.

The section closes by keeping deriving the non-Abelian gauge transformation connected to Faddeev-Popov Ghosts and BRST Symmetry.  The same general operation will be used with different objects in neighboring chapters, so the present calculation is both a result and a rehearsal.  In Faddeev-Popov Ghosts and BRST Symmetry, section 1, the useful habit is to name the object, the operation, and the assumption before using the compact notation.



\paragraph{Step-by-step derivation.}

For Faddeev-Popov Ghosts and BRST Symmetry: The concrete problem: deriving the non-Abelian gauge transformation, the derivation is written from the smallest usable model upward.

Let \(D_\mu=\p_\mu+\ii gA_\mu\), with \(A_\mu=A_\mu^aT^a\).  Acting on a test field,
\[
  [D_\mu,D_\nu]
  =
  \ii g(\p_\mu A_\nu-\p_\nu A_\mu)-g^2[A_\mu,A_\nu].
\]
Write this as \([D_\mu,D_\nu]=\ii gF_{\mu\nu}\).  The commutator term is the new feature:
\[
  F_{\mu\nu}=\p_\mu A_\nu-\p_\nu A_\mu+\ii g[A_\mu,A_\nu].
\]
Using \([T^a,T^b]=\ii f^{abc}T^c\) gives the component expression with the \(gf^{abc}A_\mu^bA_\nu^c\) self-interaction.  Squaring \(F_{\mu\nu}\) in the action produces cubic and quartic gauge-boson vertices.  In Faddeev-Popov Ghosts and BRST Symmetry, section 1, the useful habit is to name the object, the operation, and the assumption before using the compact notation.

The formulas to keep in view are

\[
  D_\mu=\p_\mu+\ii gA_\mu^aT^a
\]
\[
  F_{\mu\nu}^a=\p_\mu A_\nu^a-\p_\nu A_\mu^a-gf^{abc}A_\mu^bA_\nu^c
\]
\[
  \Lag_{\mathrm{YM}}=-\frac14F_{\mu\nu}^aF^{a\mu\nu}
\]

In this setting the calculation for Faddeev-Popov Ghosts and BRST Symmetry: The concrete problem: deriving the non-Abelian gauge transformation is complete when the finite model, the limiting step, and the normalization or boundary condition can all be named without looking back at the prose.




\begin{example}[A worked calculation for Faddeev-Popov Ghosts and BRST Symmetry: The concrete problem: deriving the non-Abelian gauge transformation]
Take two generators with \([T^1,T^2]=\ii T^3\).  If \(A_\mu=A_\mu^1T^1\) and \(A_\nu=A_\nu^2T^2\), then the commutator part of \(F_{\mu\nu}\) contains
\[
  \ii g[A_\mu,A_\nu]=\ii gA_\mu^1A_\nu^2[T^1,T^2]=-gA_\mu^1A_\nu^2T^3.
\]
Two gauge components have produced a third.  That is the algebraic origin of gauge-boson self-coupling.  In Faddeev-Popov Ghosts and BRST Symmetry, section 1, the useful habit is to name the object, the operation, and the assumption before using the compact notation.
\end{example}



\begin{exercise}
Set \(f^{abc}=0\) in the non-Abelian field strength.  What remains?  Use the notation of Faddeev-Popov Ghosts and BRST Symmetry: The concrete problem: deriving the non-Abelian gauge transformation.
\begin{answer}
The Abelian field strength \(\p_\mu A_\nu^a-\p_\nu A_\mu^a\) for each component.
\end{answer}
\end{exercise}

\begin{exercise}
Why do ghosts interact in Yang-Mills theory?  Use the notation of Faddeev-Popov Ghosts and BRST Symmetry: The concrete problem: deriving the non-Abelian gauge transformation.
\begin{answer}
The Faddeev-Popov operator contains the covariant derivative, which contains \(A_\mu\).
\end{answer}
\end{exercise}



\section{Objects, notation, and units: computing a covariant commutator}

The work of this section is computing a covariant commutator.  In the chapter heading the vocabulary is ghosts, gauge volume, BRST, physical states, but the first objects on the page are more prosaic: matrix-valued gauge fields, covariant field strengths, structure constants, ghosts, BRST symmetry, and self-interactions.  The formal notation will be useful only after it has been tied to a limit, a symmetry, a normalization condition, or an integration-by-parts step.  In Faddeev-Popov Ghosts and BRST Symmetry, section 2, the useful habit is to name the object, the operation, and the assumption before using the compact notation.

The finite model is a field with several components transformed by a position-dependent unitary matrix.  In that model no mystery is available: every derivative is an ordinary derivative with respect to one listed variable, every inner product is a sum, and every boundary assumption can be written at the endpoints.  This is why the finite model is not a toy in the dismissive sense.  It is the bookkeeping device that prevents continuum notation from becoming a fog of indices.  For Faddeev-Popov Ghosts and BRST Symmetry, section 2, that bookkeeping is deliberately modest, but it is what later keeps signs, factors of \(2\pi\), and normalization constants from turning into guesses.

The central idea for this chapter is the commutator of covariant derivatives produces a field strength that contains a gauge-field commutator.  To test that idea, strip the formula down until only the operation remains.  A sign change after raising or lowering an index means the metric has entered.  A derivative moving from one factor to another means a boundary term has been handled.  A normalization constant means that some unit object, such as a delta function, basis vector, or one-particle state, has been fixed.  Read in the setting of Faddeev-Popov Ghosts and BRST Symmetry, section 2, the line is a check on the calculation rather than an invitation to memorize another rule.

For objects, notation, and units: computing a covariant commutator, the calculation should also pass a dimensional check.  The left and right sides must carry the same units; more subtly, they must transform in the same way under the symmetries already introduced.  This is not a proof, but it is a quick way to catch wrong factors of \(2\pi\), signs in the metric, missing powers of the lattice spacing, or an omitted charge.  The same step will look more abstract later; in Faddeev-Popov Ghosts and BRST Symmetry, section 2 it is kept close to the finite calculation so the limiting move remains visible.

The bridge to quantum field theory is direct.  Yang-Mills theory supplies the strong and weak gauge sectors and introduces gauge-boson self-interactions.  Later notation will carry more indices and fewer verbal reminders, so the safe habit is to keep asking which operation is being performed and which quantity is being held fixed.  At this point in Faddeev-Popov Ghosts and BRST Symmetry, section 2, it is worth asking what would fail if the boundary condition, normalization, or symmetry assumption were changed.

One warning is worth making explicit here.  Ghost fields are not optional decoration in covariant gauges; they represent the Faddeev-Popov determinant.  This warning points to the delicate step of the derivation, not to a decorative aside.  In Faddeev-Popov Ghosts and BRST Symmetry, section 2, the calculation is meant to leave a trail from an ordinary derivative, sum, matrix product, or Gaussian integral to the field-theory formula.

The section closes by keeping computing a covariant commutator connected to Faddeev-Popov Ghosts and BRST Symmetry.  The same general operation will be used with different objects in neighboring chapters, so the present calculation is both a result and a rehearsal.  In Faddeev-Popov Ghosts and BRST Symmetry, section 2, the useful habit is to name the object, the operation, and the assumption before using the compact notation.



\paragraph{Step-by-step derivation.}

For Faddeev-Popov Ghosts and BRST Symmetry: Objects, notation, and units: computing a covariant commutator, the derivation is written from the smallest usable model upward.

Let \(D_\mu=\p_\mu+\ii gA_\mu\), with \(A_\mu=A_\mu^aT^a\).  Acting on a test field,
\[
  [D_\mu,D_\nu]
  =
  \ii g(\p_\mu A_\nu-\p_\nu A_\mu)-g^2[A_\mu,A_\nu].
\]
Write this as \([D_\mu,D_\nu]=\ii gF_{\mu\nu}\).  The commutator term is the new feature:
\[
  F_{\mu\nu}=\p_\mu A_\nu-\p_\nu A_\mu+\ii g[A_\mu,A_\nu].
\]
Using \([T^a,T^b]=\ii f^{abc}T^c\) gives the component expression with the \(gf^{abc}A_\mu^bA_\nu^c\) self-interaction.  Squaring \(F_{\mu\nu}\) in the action produces cubic and quartic gauge-boson vertices.  In Faddeev-Popov Ghosts and BRST Symmetry, section 2, the useful habit is to name the object, the operation, and the assumption before using the compact notation.

The formulas to keep in view are

\[
  D_\mu=\p_\mu+\ii gA_\mu^aT^a
\]
\[
  F_{\mu\nu}^a=\p_\mu A_\nu^a-\p_\nu A_\mu^a-gf^{abc}A_\mu^bA_\nu^c
\]
\[
  \Lag_{\mathrm{YM}}=-\frac14F_{\mu\nu}^aF^{a\mu\nu}
\]

In this setting the calculation for Faddeev-Popov Ghosts and BRST Symmetry: Objects, notation, and units: computing a covariant commutator is complete when the finite model, the limiting step, and the normalization or boundary condition can all be named without looking back at the prose.




\begin{example}[A worked calculation for Faddeev-Popov Ghosts and BRST Symmetry: Objects, notation, and units: computing a covariant commutator]
Take two generators with \([T^1,T^2]=\ii T^3\).  If \(A_\mu=A_\mu^1T^1\) and \(A_\nu=A_\nu^2T^2\), then the commutator part of \(F_{\mu\nu}\) contains
\[
  \ii g[A_\mu,A_\nu]=\ii gA_\mu^1A_\nu^2[T^1,T^2]=-gA_\mu^1A_\nu^2T^3.
\]
Two gauge components have produced a third.  That is the algebraic origin of gauge-boson self-coupling.  In Faddeev-Popov Ghosts and BRST Symmetry, section 2, the useful habit is to name the object, the operation, and the assumption before using the compact notation.
\end{example}



\begin{exercise}
Set \(f^{abc}=0\) in the non-Abelian field strength.  What remains?  Use the notation of Faddeev-Popov Ghosts and BRST Symmetry: Objects, notation, and units: computing a covariant commutator.
\begin{answer}
The Abelian field strength \(\p_\mu A_\nu^a-\p_\nu A_\mu^a\) for each component.
\end{answer}
\end{exercise}

\begin{exercise}
Why do ghosts interact in Yang-Mills theory?  Use the notation of Faddeev-Popov Ghosts and BRST Symmetry: Objects, notation, and units: computing a covariant commutator.
\begin{answer}
The Faddeev-Popov operator contains the covariant derivative, which contains \(A_\mu\).
\end{answer}
\end{exercise}



\section{The finite calculation: expanding Yang-Mills interactions}

The work of this section is expanding Yang-Mills interactions.  In the chapter heading the vocabulary is ghosts, gauge volume, BRST, physical states, but the first objects on the page are more prosaic: matrix-valued gauge fields, covariant field strengths, structure constants, ghosts, BRST symmetry, and self-interactions.  The formal notation will be useful only after it has been tied to a limit, a symmetry, a normalization condition, or an integration-by-parts step.  In Faddeev-Popov Ghosts and BRST Symmetry, section 3, the useful habit is to name the object, the operation, and the assumption before using the compact notation.

The finite model is a field with several components transformed by a position-dependent unitary matrix.  In that model no mystery is available: every derivative is an ordinary derivative with respect to one listed variable, every inner product is a sum, and every boundary assumption can be written at the endpoints.  This is why the finite model is not a toy in the dismissive sense.  It is the bookkeeping device that prevents continuum notation from becoming a fog of indices.  For Faddeev-Popov Ghosts and BRST Symmetry, section 3, that bookkeeping is deliberately modest, but it is what later keeps signs, factors of \(2\pi\), and normalization constants from turning into guesses.

The central idea for this chapter is the commutator of covariant derivatives produces a field strength that contains a gauge-field commutator.  To test that idea, strip the formula down until only the operation remains.  A sign change after raising or lowering an index means the metric has entered.  A derivative moving from one factor to another means a boundary term has been handled.  A normalization constant means that some unit object, such as a delta function, basis vector, or one-particle state, has been fixed.  Read in the setting of Faddeev-Popov Ghosts and BRST Symmetry, section 3, the line is a check on the calculation rather than an invitation to memorize another rule.

For the finite calculation: expanding yang-mills interactions, the calculation should also pass a dimensional check.  The left and right sides must carry the same units; more subtly, they must transform in the same way under the symmetries already introduced.  This is not a proof, but it is a quick way to catch wrong factors of \(2\pi\), signs in the metric, missing powers of the lattice spacing, or an omitted charge.  The same step will look more abstract later; in Faddeev-Popov Ghosts and BRST Symmetry, section 3 it is kept close to the finite calculation so the limiting move remains visible.

The bridge to quantum field theory is direct.  Yang-Mills theory supplies the strong and weak gauge sectors and introduces gauge-boson self-interactions.  Later notation will carry more indices and fewer verbal reminders, so the safe habit is to keep asking which operation is being performed and which quantity is being held fixed.  At this point in Faddeev-Popov Ghosts and BRST Symmetry, section 3, it is worth asking what would fail if the boundary condition, normalization, or symmetry assumption were changed.

One warning is worth making explicit here.  Ghost fields are not optional decoration in covariant gauges; they represent the Faddeev-Popov determinant.  This warning points to the delicate step of the derivation, not to a decorative aside.  In Faddeev-Popov Ghosts and BRST Symmetry, section 3, the calculation is meant to leave a trail from an ordinary derivative, sum, matrix product, or Gaussian integral to the field-theory formula.

The section closes by keeping expanding Yang-Mills interactions connected to Faddeev-Popov Ghosts and BRST Symmetry.  The same general operation will be used with different objects in neighboring chapters, so the present calculation is both a result and a rehearsal.  In Faddeev-Popov Ghosts and BRST Symmetry, section 3, the useful habit is to name the object, the operation, and the assumption before using the compact notation.



\paragraph{Step-by-step derivation.}

For Faddeev-Popov Ghosts and BRST Symmetry: The finite calculation: expanding Yang-Mills interactions, the derivation is written from the smallest usable model upward.

Let \(D_\mu=\p_\mu+\ii gA_\mu\), with \(A_\mu=A_\mu^aT^a\).  Acting on a test field,
\[
  [D_\mu,D_\nu]
  =
  \ii g(\p_\mu A_\nu-\p_\nu A_\mu)-g^2[A_\mu,A_\nu].
\]
Write this as \([D_\mu,D_\nu]=\ii gF_{\mu\nu}\).  The commutator term is the new feature:
\[
  F_{\mu\nu}=\p_\mu A_\nu-\p_\nu A_\mu+\ii g[A_\mu,A_\nu].
\]
Using \([T^a,T^b]=\ii f^{abc}T^c\) gives the component expression with the \(gf^{abc}A_\mu^bA_\nu^c\) self-interaction.  Squaring \(F_{\mu\nu}\) in the action produces cubic and quartic gauge-boson vertices.  In Faddeev-Popov Ghosts and BRST Symmetry, section 3, the useful habit is to name the object, the operation, and the assumption before using the compact notation.

The formulas to keep in view are

\[
  D_\mu=\p_\mu+\ii gA_\mu^aT^a
\]
\[
  F_{\mu\nu}^a=\p_\mu A_\nu^a-\p_\nu A_\mu^a-gf^{abc}A_\mu^bA_\nu^c
\]
\[
  \Lag_{\mathrm{YM}}=-\frac14F_{\mu\nu}^aF^{a\mu\nu}
\]

In this setting the calculation for Faddeev-Popov Ghosts and BRST Symmetry: The finite calculation: expanding Yang-Mills interactions is complete when the finite model, the limiting step, and the normalization or boundary condition can all be named without looking back at the prose.




\begin{example}[A worked calculation for Faddeev-Popov Ghosts and BRST Symmetry: The finite calculation: expanding Yang-Mills interactions]
Take two generators with \([T^1,T^2]=\ii T^3\).  If \(A_\mu=A_\mu^1T^1\) and \(A_\nu=A_\nu^2T^2\), then the commutator part of \(F_{\mu\nu}\) contains
\[
  \ii g[A_\mu,A_\nu]=\ii gA_\mu^1A_\nu^2[T^1,T^2]=-gA_\mu^1A_\nu^2T^3.
\]
Two gauge components have produced a third.  That is the algebraic origin of gauge-boson self-coupling.  In Faddeev-Popov Ghosts and BRST Symmetry, section 3, the useful habit is to name the object, the operation, and the assumption before using the compact notation.
\end{example}



\begin{exercise}
Set \(f^{abc}=0\) in the non-Abelian field strength.  What remains?  Use the notation of Faddeev-Popov Ghosts and BRST Symmetry: The finite calculation: expanding Yang-Mills interactions.
\begin{answer}
The Abelian field strength \(\p_\mu A_\nu^a-\p_\nu A_\mu^a\) for each component.
\end{answer}
\end{exercise}

\begin{exercise}
Why do ghosts interact in Yang-Mills theory?  Use the notation of Faddeev-Popov Ghosts and BRST Symmetry: The finite calculation: expanding Yang-Mills interactions.
\begin{answer}
The Faddeev-Popov operator contains the covariant derivative, which contains \(A_\mu\).
\end{answer}
\end{exercise}



\section{The central derivation: varying the Yang-Mills action}

The work of this section is varying the Yang-Mills action.  In the chapter heading the vocabulary is ghosts, gauge volume, BRST, physical states, but the first objects on the page are more prosaic: matrix-valued gauge fields, covariant field strengths, structure constants, ghosts, BRST symmetry, and self-interactions.  The formal notation will be useful only after it has been tied to a limit, a symmetry, a normalization condition, or an integration-by-parts step.  In Faddeev-Popov Ghosts and BRST Symmetry, section 4, the useful habit is to name the object, the operation, and the assumption before using the compact notation.

The finite model is a field with several components transformed by a position-dependent unitary matrix.  In that model no mystery is available: every derivative is an ordinary derivative with respect to one listed variable, every inner product is a sum, and every boundary assumption can be written at the endpoints.  This is why the finite model is not a toy in the dismissive sense.  It is the bookkeeping device that prevents continuum notation from becoming a fog of indices.  For Faddeev-Popov Ghosts and BRST Symmetry, section 4, that bookkeeping is deliberately modest, but it is what later keeps signs, factors of \(2\pi\), and normalization constants from turning into guesses.

The central idea for this chapter is the commutator of covariant derivatives produces a field strength that contains a gauge-field commutator.  To test that idea, strip the formula down until only the operation remains.  A sign change after raising or lowering an index means the metric has entered.  A derivative moving from one factor to another means a boundary term has been handled.  A normalization constant means that some unit object, such as a delta function, basis vector, or one-particle state, has been fixed.  Read in the setting of Faddeev-Popov Ghosts and BRST Symmetry, section 4, the line is a check on the calculation rather than an invitation to memorize another rule.

For the central derivation: varying the yang-mills action, the calculation should also pass a dimensional check.  The left and right sides must carry the same units; more subtly, they must transform in the same way under the symmetries already introduced.  This is not a proof, but it is a quick way to catch wrong factors of \(2\pi\), signs in the metric, missing powers of the lattice spacing, or an omitted charge.  The same step will look more abstract later; in Faddeev-Popov Ghosts and BRST Symmetry, section 4 it is kept close to the finite calculation so the limiting move remains visible.

The bridge to quantum field theory is direct.  Yang-Mills theory supplies the strong and weak gauge sectors and introduces gauge-boson self-interactions.  Later notation will carry more indices and fewer verbal reminders, so the safe habit is to keep asking which operation is being performed and which quantity is being held fixed.  At this point in Faddeev-Popov Ghosts and BRST Symmetry, section 4, it is worth asking what would fail if the boundary condition, normalization, or symmetry assumption were changed.

One warning is worth making explicit here.  Ghost fields are not optional decoration in covariant gauges; they represent the Faddeev-Popov determinant.  This warning points to the delicate step of the derivation, not to a decorative aside.  In Faddeev-Popov Ghosts and BRST Symmetry, section 4, the calculation is meant to leave a trail from an ordinary derivative, sum, matrix product, or Gaussian integral to the field-theory formula.

The section closes by keeping varying the Yang-Mills action connected to Faddeev-Popov Ghosts and BRST Symmetry.  The same general operation will be used with different objects in neighboring chapters, so the present calculation is both a result and a rehearsal.  In Faddeev-Popov Ghosts and BRST Symmetry, section 4, the useful habit is to name the object, the operation, and the assumption before using the compact notation.



\paragraph{Step-by-step derivation.}

For Faddeev-Popov Ghosts and BRST Symmetry: The central derivation: varying the Yang-Mills action, the derivation is written from the smallest usable model upward.

Let \(D_\mu=\p_\mu+\ii gA_\mu\), with \(A_\mu=A_\mu^aT^a\).  Acting on a test field,
\[
  [D_\mu,D_\nu]
  =
  \ii g(\p_\mu A_\nu-\p_\nu A_\mu)-g^2[A_\mu,A_\nu].
\]
Write this as \([D_\mu,D_\nu]=\ii gF_{\mu\nu}\).  The commutator term is the new feature:
\[
  F_{\mu\nu}=\p_\mu A_\nu-\p_\nu A_\mu+\ii g[A_\mu,A_\nu].
\]
Using \([T^a,T^b]=\ii f^{abc}T^c\) gives the component expression with the \(gf^{abc}A_\mu^bA_\nu^c\) self-interaction.  Squaring \(F_{\mu\nu}\) in the action produces cubic and quartic gauge-boson vertices.  In Faddeev-Popov Ghosts and BRST Symmetry, section 4, the useful habit is to name the object, the operation, and the assumption before using the compact notation.

The formulas to keep in view are

\[
  D_\mu=\p_\mu+\ii gA_\mu^aT^a
\]
\[
  F_{\mu\nu}^a=\p_\mu A_\nu^a-\p_\nu A_\mu^a-gf^{abc}A_\mu^bA_\nu^c
\]
\[
  \Lag_{\mathrm{YM}}=-\frac14F_{\mu\nu}^aF^{a\mu\nu}
\]

In this setting the calculation for Faddeev-Popov Ghosts and BRST Symmetry: The central derivation: varying the Yang-Mills action is complete when the finite model, the limiting step, and the normalization or boundary condition can all be named without looking back at the prose.




\begin{example}[A worked calculation for Faddeev-Popov Ghosts and BRST Symmetry: The central derivation: varying the Yang-Mills action]
Take two generators with \([T^1,T^2]=\ii T^3\).  If \(A_\mu=A_\mu^1T^1\) and \(A_\nu=A_\nu^2T^2\), then the commutator part of \(F_{\mu\nu}\) contains
\[
  \ii g[A_\mu,A_\nu]=\ii gA_\mu^1A_\nu^2[T^1,T^2]=-gA_\mu^1A_\nu^2T^3.
\]
Two gauge components have produced a third.  That is the algebraic origin of gauge-boson self-coupling.  In Faddeev-Popov Ghosts and BRST Symmetry, section 4, the useful habit is to name the object, the operation, and the assumption before using the compact notation.
\end{example}



\begin{exercise}
Set \(f^{abc}=0\) in the non-Abelian field strength.  What remains?  Use the notation of Faddeev-Popov Ghosts and BRST Symmetry: The central derivation: varying the Yang-Mills action.
\begin{answer}
The Abelian field strength \(\p_\mu A_\nu^a-\p_\nu A_\mu^a\) for each component.
\end{answer}
\end{exercise}

\begin{exercise}
Why do ghosts interact in Yang-Mills theory?  Use the notation of Faddeev-Popov Ghosts and BRST Symmetry: The central derivation: varying the Yang-Mills action.
\begin{answer}
The Faddeev-Popov operator contains the covariant derivative, which contains \(A_\mu\).
\end{answer}
\end{exercise}



\section{Worked calculation: fixing a non-Abelian gauge}

The work of this section is fixing a non-Abelian gauge.  In the chapter heading the vocabulary is ghosts, gauge volume, BRST, physical states, but the first objects on the page are more prosaic: matrix-valued gauge fields, covariant field strengths, structure constants, ghosts, BRST symmetry, and self-interactions.  The formal notation will be useful only after it has been tied to a limit, a symmetry, a normalization condition, or an integration-by-parts step.  In Faddeev-Popov Ghosts and BRST Symmetry, section 5, the useful habit is to name the object, the operation, and the assumption before using the compact notation.

The finite model is a field with several components transformed by a position-dependent unitary matrix.  In that model no mystery is available: every derivative is an ordinary derivative with respect to one listed variable, every inner product is a sum, and every boundary assumption can be written at the endpoints.  This is why the finite model is not a toy in the dismissive sense.  It is the bookkeeping device that prevents continuum notation from becoming a fog of indices.  For Faddeev-Popov Ghosts and BRST Symmetry, section 5, that bookkeeping is deliberately modest, but it is what later keeps signs, factors of \(2\pi\), and normalization constants from turning into guesses.

The central idea for this chapter is the commutator of covariant derivatives produces a field strength that contains a gauge-field commutator.  To test that idea, strip the formula down until only the operation remains.  A sign change after raising or lowering an index means the metric has entered.  A derivative moving from one factor to another means a boundary term has been handled.  A normalization constant means that some unit object, such as a delta function, basis vector, or one-particle state, has been fixed.  Read in the setting of Faddeev-Popov Ghosts and BRST Symmetry, section 5, the line is a check on the calculation rather than an invitation to memorize another rule.

For worked calculation: fixing a non-abelian gauge, the calculation should also pass a dimensional check.  The left and right sides must carry the same units; more subtly, they must transform in the same way under the symmetries already introduced.  This is not a proof, but it is a quick way to catch wrong factors of \(2\pi\), signs in the metric, missing powers of the lattice spacing, or an omitted charge.  The same step will look more abstract later; in Faddeev-Popov Ghosts and BRST Symmetry, section 5 it is kept close to the finite calculation so the limiting move remains visible.

The bridge to quantum field theory is direct.  Yang-Mills theory supplies the strong and weak gauge sectors and introduces gauge-boson self-interactions.  Later notation will carry more indices and fewer verbal reminders, so the safe habit is to keep asking which operation is being performed and which quantity is being held fixed.  At this point in Faddeev-Popov Ghosts and BRST Symmetry, section 5, it is worth asking what would fail if the boundary condition, normalization, or symmetry assumption were changed.

One warning is worth making explicit here.  Ghost fields are not optional decoration in covariant gauges; they represent the Faddeev-Popov determinant.  This warning points to the delicate step of the derivation, not to a decorative aside.  In Faddeev-Popov Ghosts and BRST Symmetry, section 5, the calculation is meant to leave a trail from an ordinary derivative, sum, matrix product, or Gaussian integral to the field-theory formula.

The section closes by keeping fixing a non-Abelian gauge connected to Faddeev-Popov Ghosts and BRST Symmetry.  The same general operation will be used with different objects in neighboring chapters, so the present calculation is both a result and a rehearsal.  In Faddeev-Popov Ghosts and BRST Symmetry, section 5, the useful habit is to name the object, the operation, and the assumption before using the compact notation.



\paragraph{Step-by-step derivation.}

For Faddeev-Popov Ghosts and BRST Symmetry: Worked calculation: fixing a non-Abelian gauge, the derivation is written from the smallest usable model upward.

Let \(D_\mu=\p_\mu+\ii gA_\mu\), with \(A_\mu=A_\mu^aT^a\).  Acting on a test field,
\[
  [D_\mu,D_\nu]
  =
  \ii g(\p_\mu A_\nu-\p_\nu A_\mu)-g^2[A_\mu,A_\nu].
\]
Write this as \([D_\mu,D_\nu]=\ii gF_{\mu\nu}\).  The commutator term is the new feature:
\[
  F_{\mu\nu}=\p_\mu A_\nu-\p_\nu A_\mu+\ii g[A_\mu,A_\nu].
\]
Using \([T^a,T^b]=\ii f^{abc}T^c\) gives the component expression with the \(gf^{abc}A_\mu^bA_\nu^c\) self-interaction.  Squaring \(F_{\mu\nu}\) in the action produces cubic and quartic gauge-boson vertices.  In Faddeev-Popov Ghosts and BRST Symmetry, section 5, the useful habit is to name the object, the operation, and the assumption before using the compact notation.

The formulas to keep in view are

\[
  D_\mu=\p_\mu+\ii gA_\mu^aT^a
\]
\[
  F_{\mu\nu}^a=\p_\mu A_\nu^a-\p_\nu A_\mu^a-gf^{abc}A_\mu^bA_\nu^c
\]
\[
  \Lag_{\mathrm{YM}}=-\frac14F_{\mu\nu}^aF^{a\mu\nu}
\]

In this setting the calculation for Faddeev-Popov Ghosts and BRST Symmetry: Worked calculation: fixing a non-Abelian gauge is complete when the finite model, the limiting step, and the normalization or boundary condition can all be named without looking back at the prose.




\begin{example}[A worked calculation for Faddeev-Popov Ghosts and BRST Symmetry: Worked calculation: fixing a non-Abelian gauge]
Take two generators with \([T^1,T^2]=\ii T^3\).  If \(A_\mu=A_\mu^1T^1\) and \(A_\nu=A_\nu^2T^2\), then the commutator part of \(F_{\mu\nu}\) contains
\[
  \ii g[A_\mu,A_\nu]=\ii gA_\mu^1A_\nu^2[T^1,T^2]=-gA_\mu^1A_\nu^2T^3.
\]
Two gauge components have produced a third.  That is the algebraic origin of gauge-boson self-coupling.  In Faddeev-Popov Ghosts and BRST Symmetry, section 5, the useful habit is to name the object, the operation, and the assumption before using the compact notation.
\end{example}



\begin{exercise}
Set \(f^{abc}=0\) in the non-Abelian field strength.  What remains?  Use the notation of Faddeev-Popov Ghosts and BRST Symmetry: Worked calculation: fixing a non-Abelian gauge.
\begin{answer}
The Abelian field strength \(\p_\mu A_\nu^a-\p_\nu A_\mu^a\) for each component.
\end{answer}
\end{exercise}

\begin{exercise}
Why do ghosts interact in Yang-Mills theory?  Use the notation of Faddeev-Popov Ghosts and BRST Symmetry: Worked calculation: fixing a non-Abelian gauge.
\begin{answer}
The Faddeev-Popov operator contains the covariant derivative, which contains \(A_\mu\).
\end{answer}
\end{exercise}



\section{Boundary conditions, normalization, and signs: introducing ghosts}

The work of this section is introducing ghosts.  In the chapter heading the vocabulary is ghosts, gauge volume, BRST, physical states, but the first objects on the page are more prosaic: matrix-valued gauge fields, covariant field strengths, structure constants, ghosts, BRST symmetry, and self-interactions.  The formal notation will be useful only after it has been tied to a limit, a symmetry, a normalization condition, or an integration-by-parts step.  In Faddeev-Popov Ghosts and BRST Symmetry, section 6, the useful habit is to name the object, the operation, and the assumption before using the compact notation.

The finite model is a field with several components transformed by a position-dependent unitary matrix.  In that model no mystery is available: every derivative is an ordinary derivative with respect to one listed variable, every inner product is a sum, and every boundary assumption can be written at the endpoints.  This is why the finite model is not a toy in the dismissive sense.  It is the bookkeeping device that prevents continuum notation from becoming a fog of indices.  For Faddeev-Popov Ghosts and BRST Symmetry, section 6, that bookkeeping is deliberately modest, but it is what later keeps signs, factors of \(2\pi\), and normalization constants from turning into guesses.

The central idea for this chapter is the commutator of covariant derivatives produces a field strength that contains a gauge-field commutator.  To test that idea, strip the formula down until only the operation remains.  A sign change after raising or lowering an index means the metric has entered.  A derivative moving from one factor to another means a boundary term has been handled.  A normalization constant means that some unit object, such as a delta function, basis vector, or one-particle state, has been fixed.  Read in the setting of Faddeev-Popov Ghosts and BRST Symmetry, section 6, the line is a check on the calculation rather than an invitation to memorize another rule.

For boundary conditions, normalization, and signs: introducing ghosts, the calculation should also pass a dimensional check.  The left and right sides must carry the same units; more subtly, they must transform in the same way under the symmetries already introduced.  This is not a proof, but it is a quick way to catch wrong factors of \(2\pi\), signs in the metric, missing powers of the lattice spacing, or an omitted charge.  The same step will look more abstract later; in Faddeev-Popov Ghosts and BRST Symmetry, section 6 it is kept close to the finite calculation so the limiting move remains visible.

The bridge to quantum field theory is direct.  Yang-Mills theory supplies the strong and weak gauge sectors and introduces gauge-boson self-interactions.  Later notation will carry more indices and fewer verbal reminders, so the safe habit is to keep asking which operation is being performed and which quantity is being held fixed.  At this point in Faddeev-Popov Ghosts and BRST Symmetry, section 6, it is worth asking what would fail if the boundary condition, normalization, or symmetry assumption were changed.

One warning is worth making explicit here.  Ghost fields are not optional decoration in covariant gauges; they represent the Faddeev-Popov determinant.  This warning points to the delicate step of the derivation, not to a decorative aside.  In Faddeev-Popov Ghosts and BRST Symmetry, section 6, the calculation is meant to leave a trail from an ordinary derivative, sum, matrix product, or Gaussian integral to the field-theory formula.

The section closes by keeping introducing ghosts connected to Faddeev-Popov Ghosts and BRST Symmetry.  The same general operation will be used with different objects in neighboring chapters, so the present calculation is both a result and a rehearsal.  In Faddeev-Popov Ghosts and BRST Symmetry, section 6, the useful habit is to name the object, the operation, and the assumption before using the compact notation.



\paragraph{Step-by-step derivation.}

For Faddeev-Popov Ghosts and BRST Symmetry: Boundary conditions, normalization, and signs: introducing ghosts, the derivation is written from the smallest usable model upward.

Let \(D_\mu=\p_\mu+\ii gA_\mu\), with \(A_\mu=A_\mu^aT^a\).  Acting on a test field,
\[
  [D_\mu,D_\nu]
  =
  \ii g(\p_\mu A_\nu-\p_\nu A_\mu)-g^2[A_\mu,A_\nu].
\]
Write this as \([D_\mu,D_\nu]=\ii gF_{\mu\nu}\).  The commutator term is the new feature:
\[
  F_{\mu\nu}=\p_\mu A_\nu-\p_\nu A_\mu+\ii g[A_\mu,A_\nu].
\]
Using \([T^a,T^b]=\ii f^{abc}T^c\) gives the component expression with the \(gf^{abc}A_\mu^bA_\nu^c\) self-interaction.  Squaring \(F_{\mu\nu}\) in the action produces cubic and quartic gauge-boson vertices.  In Faddeev-Popov Ghosts and BRST Symmetry, section 6, the useful habit is to name the object, the operation, and the assumption before using the compact notation.

The formulas to keep in view are

\[
  D_\mu=\p_\mu+\ii gA_\mu^aT^a
\]
\[
  F_{\mu\nu}^a=\p_\mu A_\nu^a-\p_\nu A_\mu^a-gf^{abc}A_\mu^bA_\nu^c
\]
\[
  \Lag_{\mathrm{YM}}=-\frac14F_{\mu\nu}^aF^{a\mu\nu}
\]

In this setting the calculation for Faddeev-Popov Ghosts and BRST Symmetry: Boundary conditions, normalization, and signs: introducing ghosts is complete when the finite model, the limiting step, and the normalization or boundary condition can all be named without looking back at the prose.




\begin{example}[A worked calculation for Faddeev-Popov Ghosts and BRST Symmetry: Boundary conditions, normalization, and signs: introducing ghosts]
Take two generators with \([T^1,T^2]=\ii T^3\).  If \(A_\mu=A_\mu^1T^1\) and \(A_\nu=A_\nu^2T^2\), then the commutator part of \(F_{\mu\nu}\) contains
\[
  \ii g[A_\mu,A_\nu]=\ii gA_\mu^1A_\nu^2[T^1,T^2]=-gA_\mu^1A_\nu^2T^3.
\]
Two gauge components have produced a third.  That is the algebraic origin of gauge-boson self-coupling.  In Faddeev-Popov Ghosts and BRST Symmetry, section 6, the useful habit is to name the object, the operation, and the assumption before using the compact notation.
\end{example}



\begin{exercise}
Set \(f^{abc}=0\) in the non-Abelian field strength.  What remains?  Use the notation of Faddeev-Popov Ghosts and BRST Symmetry: Boundary conditions, normalization, and signs: introducing ghosts.
\begin{answer}
The Abelian field strength \(\p_\mu A_\nu^a-\p_\nu A_\mu^a\) for each component.
\end{answer}
\end{exercise}

\begin{exercise}
Why do ghosts interact in Yang-Mills theory?  Use the notation of Faddeev-Popov Ghosts and BRST Symmetry: Boundary conditions, normalization, and signs: introducing ghosts.
\begin{answer}
The Faddeev-Popov operator contains the covariant derivative, which contains \(A_\mu\).
\end{answer}
\end{exercise}



\section{How the idea enters quantum field theory: checking BRST nilpotency}

The work of this section is checking BRST nilpotency.  In the chapter heading the vocabulary is ghosts, gauge volume, BRST, physical states, but the first objects on the page are more prosaic: matrix-valued gauge fields, covariant field strengths, structure constants, ghosts, BRST symmetry, and self-interactions.  The formal notation will be useful only after it has been tied to a limit, a symmetry, a normalization condition, or an integration-by-parts step.  In Faddeev-Popov Ghosts and BRST Symmetry, section 7, the useful habit is to name the object, the operation, and the assumption before using the compact notation.

The finite model is a field with several components transformed by a position-dependent unitary matrix.  In that model no mystery is available: every derivative is an ordinary derivative with respect to one listed variable, every inner product is a sum, and every boundary assumption can be written at the endpoints.  This is why the finite model is not a toy in the dismissive sense.  It is the bookkeeping device that prevents continuum notation from becoming a fog of indices.  For Faddeev-Popov Ghosts and BRST Symmetry, section 7, that bookkeeping is deliberately modest, but it is what later keeps signs, factors of \(2\pi\), and normalization constants from turning into guesses.

The central idea for this chapter is the commutator of covariant derivatives produces a field strength that contains a gauge-field commutator.  To test that idea, strip the formula down until only the operation remains.  A sign change after raising or lowering an index means the metric has entered.  A derivative moving from one factor to another means a boundary term has been handled.  A normalization constant means that some unit object, such as a delta function, basis vector, or one-particle state, has been fixed.  Read in the setting of Faddeev-Popov Ghosts and BRST Symmetry, section 7, the line is a check on the calculation rather than an invitation to memorize another rule.

For how the idea enters quantum field theory: checking brst nilpotency, the calculation should also pass a dimensional check.  The left and right sides must carry the same units; more subtly, they must transform in the same way under the symmetries already introduced.  This is not a proof, but it is a quick way to catch wrong factors of \(2\pi\), signs in the metric, missing powers of the lattice spacing, or an omitted charge.  The same step will look more abstract later; in Faddeev-Popov Ghosts and BRST Symmetry, section 7 it is kept close to the finite calculation so the limiting move remains visible.

The bridge to quantum field theory is direct.  Yang-Mills theory supplies the strong and weak gauge sectors and introduces gauge-boson self-interactions.  Later notation will carry more indices and fewer verbal reminders, so the safe habit is to keep asking which operation is being performed and which quantity is being held fixed.  At this point in Faddeev-Popov Ghosts and BRST Symmetry, section 7, it is worth asking what would fail if the boundary condition, normalization, or symmetry assumption were changed.

One warning is worth making explicit here.  Ghost fields are not optional decoration in covariant gauges; they represent the Faddeev-Popov determinant.  This warning points to the delicate step of the derivation, not to a decorative aside.  In Faddeev-Popov Ghosts and BRST Symmetry, section 7, the calculation is meant to leave a trail from an ordinary derivative, sum, matrix product, or Gaussian integral to the field-theory formula.

The section closes by keeping checking BRST nilpotency connected to Faddeev-Popov Ghosts and BRST Symmetry.  The same general operation will be used with different objects in neighboring chapters, so the present calculation is both a result and a rehearsal.  In Faddeev-Popov Ghosts and BRST Symmetry, section 7, the useful habit is to name the object, the operation, and the assumption before using the compact notation.



\paragraph{Step-by-step derivation.}

For Faddeev-Popov Ghosts and BRST Symmetry: How the idea enters quantum field theory: checking BRST nilpotency, the derivation is written from the smallest usable model upward.

Let \(D_\mu=\p_\mu+\ii gA_\mu\), with \(A_\mu=A_\mu^aT^a\).  Acting on a test field,
\[
  [D_\mu,D_\nu]
  =
  \ii g(\p_\mu A_\nu-\p_\nu A_\mu)-g^2[A_\mu,A_\nu].
\]
Write this as \([D_\mu,D_\nu]=\ii gF_{\mu\nu}\).  The commutator term is the new feature:
\[
  F_{\mu\nu}=\p_\mu A_\nu-\p_\nu A_\mu+\ii g[A_\mu,A_\nu].
\]
Using \([T^a,T^b]=\ii f^{abc}T^c\) gives the component expression with the \(gf^{abc}A_\mu^bA_\nu^c\) self-interaction.  Squaring \(F_{\mu\nu}\) in the action produces cubic and quartic gauge-boson vertices.  In Faddeev-Popov Ghosts and BRST Symmetry, section 7, the useful habit is to name the object, the operation, and the assumption before using the compact notation.

The formulas to keep in view are

\[
  D_\mu=\p_\mu+\ii gA_\mu^aT^a
\]
\[
  F_{\mu\nu}^a=\p_\mu A_\nu^a-\p_\nu A_\mu^a-gf^{abc}A_\mu^bA_\nu^c
\]
\[
  \Lag_{\mathrm{YM}}=-\frac14F_{\mu\nu}^aF^{a\mu\nu}
\]

In this setting the calculation for Faddeev-Popov Ghosts and BRST Symmetry: How the idea enters quantum field theory: checking BRST nilpotency is complete when the finite model, the limiting step, and the normalization or boundary condition can all be named without looking back at the prose.




\begin{example}[A worked calculation for Faddeev-Popov Ghosts and BRST Symmetry: How the idea enters quantum field theory: checking BRST nilpotency]
Take two generators with \([T^1,T^2]=\ii T^3\).  If \(A_\mu=A_\mu^1T^1\) and \(A_\nu=A_\nu^2T^2\), then the commutator part of \(F_{\mu\nu}\) contains
\[
  \ii g[A_\mu,A_\nu]=\ii gA_\mu^1A_\nu^2[T^1,T^2]=-gA_\mu^1A_\nu^2T^3.
\]
Two gauge components have produced a third.  That is the algebraic origin of gauge-boson self-coupling.  In Faddeev-Popov Ghosts and BRST Symmetry, section 7, the useful habit is to name the object, the operation, and the assumption before using the compact notation.
\end{example}



\begin{exercise}
Set \(f^{abc}=0\) in the non-Abelian field strength.  What remains?  Use the notation of Faddeev-Popov Ghosts and BRST Symmetry: How the idea enters quantum field theory: checking BRST nilpotency.
\begin{answer}
The Abelian field strength \(\p_\mu A_\nu^a-\p_\nu A_\mu^a\) for each component.
\end{answer}
\end{exercise}

\begin{exercise}
Why do ghosts interact in Yang-Mills theory?  Use the notation of Faddeev-Popov Ghosts and BRST Symmetry: How the idea enters quantum field theory: checking BRST nilpotency.
\begin{answer}
The Faddeev-Popov operator contains the covariant derivative, which contains \(A_\mu\).
\end{answer}
\end{exercise}



\section{Review problems and extensions: reading asymptotic freedom}

The work of this section is reading asymptotic freedom.  In the chapter heading the vocabulary is ghosts, gauge volume, BRST, physical states, but the first objects on the page are more prosaic: matrix-valued gauge fields, covariant field strengths, structure constants, ghosts, BRST symmetry, and self-interactions.  The formal notation will be useful only after it has been tied to a limit, a symmetry, a normalization condition, or an integration-by-parts step.  In Faddeev-Popov Ghosts and BRST Symmetry, section 8, the useful habit is to name the object, the operation, and the assumption before using the compact notation.

The finite model is a field with several components transformed by a position-dependent unitary matrix.  In that model no mystery is available: every derivative is an ordinary derivative with respect to one listed variable, every inner product is a sum, and every boundary assumption can be written at the endpoints.  This is why the finite model is not a toy in the dismissive sense.  It is the bookkeeping device that prevents continuum notation from becoming a fog of indices.  For Faddeev-Popov Ghosts and BRST Symmetry, section 8, that bookkeeping is deliberately modest, but it is what later keeps signs, factors of \(2\pi\), and normalization constants from turning into guesses.

The central idea for this chapter is the commutator of covariant derivatives produces a field strength that contains a gauge-field commutator.  To test that idea, strip the formula down until only the operation remains.  A sign change after raising or lowering an index means the metric has entered.  A derivative moving from one factor to another means a boundary term has been handled.  A normalization constant means that some unit object, such as a delta function, basis vector, or one-particle state, has been fixed.  Read in the setting of Faddeev-Popov Ghosts and BRST Symmetry, section 8, the line is a check on the calculation rather than an invitation to memorize another rule.

For review problems and extensions: reading asymptotic freedom, the calculation should also pass a dimensional check.  The left and right sides must carry the same units; more subtly, they must transform in the same way under the symmetries already introduced.  This is not a proof, but it is a quick way to catch wrong factors of \(2\pi\), signs in the metric, missing powers of the lattice spacing, or an omitted charge.  The same step will look more abstract later; in Faddeev-Popov Ghosts and BRST Symmetry, section 8 it is kept close to the finite calculation so the limiting move remains visible.

The bridge to quantum field theory is direct.  Yang-Mills theory supplies the strong and weak gauge sectors and introduces gauge-boson self-interactions.  Later notation will carry more indices and fewer verbal reminders, so the safe habit is to keep asking which operation is being performed and which quantity is being held fixed.  At this point in Faddeev-Popov Ghosts and BRST Symmetry, section 8, it is worth asking what would fail if the boundary condition, normalization, or symmetry assumption were changed.

One warning is worth making explicit here.  Ghost fields are not optional decoration in covariant gauges; they represent the Faddeev-Popov determinant.  This warning points to the delicate step of the derivation, not to a decorative aside.  In Faddeev-Popov Ghosts and BRST Symmetry, section 8, the calculation is meant to leave a trail from an ordinary derivative, sum, matrix product, or Gaussian integral to the field-theory formula.

The section closes by keeping reading asymptotic freedom connected to Faddeev-Popov Ghosts and BRST Symmetry.  The same general operation will be used with different objects in neighboring chapters, so the present calculation is both a result and a rehearsal.  In Faddeev-Popov Ghosts and BRST Symmetry, section 8, the useful habit is to name the object, the operation, and the assumption before using the compact notation.



\paragraph{Step-by-step derivation.}

For Faddeev-Popov Ghosts and BRST Symmetry: Review problems and extensions: reading asymptotic freedom, the derivation is written from the smallest usable model upward.

Let \(D_\mu=\p_\mu+\ii gA_\mu\), with \(A_\mu=A_\mu^aT^a\).  Acting on a test field,
\[
  [D_\mu,D_\nu]
  =
  \ii g(\p_\mu A_\nu-\p_\nu A_\mu)-g^2[A_\mu,A_\nu].
\]
Write this as \([D_\mu,D_\nu]=\ii gF_{\mu\nu}\).  The commutator term is the new feature:
\[
  F_{\mu\nu}=\p_\mu A_\nu-\p_\nu A_\mu+\ii g[A_\mu,A_\nu].
\]
Using \([T^a,T^b]=\ii f^{abc}T^c\) gives the component expression with the \(gf^{abc}A_\mu^bA_\nu^c\) self-interaction.  Squaring \(F_{\mu\nu}\) in the action produces cubic and quartic gauge-boson vertices.  In Faddeev-Popov Ghosts and BRST Symmetry, section 8, the useful habit is to name the object, the operation, and the assumption before using the compact notation.

The formulas to keep in view are

\[
  D_\mu=\p_\mu+\ii gA_\mu^aT^a
\]
\[
  F_{\mu\nu}^a=\p_\mu A_\nu^a-\p_\nu A_\mu^a-gf^{abc}A_\mu^bA_\nu^c
\]
\[
  \Lag_{\mathrm{YM}}=-\frac14F_{\mu\nu}^aF^{a\mu\nu}
\]

In this setting the calculation for Faddeev-Popov Ghosts and BRST Symmetry: Review problems and extensions: reading asymptotic freedom is complete when the finite model, the limiting step, and the normalization or boundary condition can all be named without looking back at the prose.




\begin{example}[A worked calculation for Faddeev-Popov Ghosts and BRST Symmetry: Review problems and extensions: reading asymptotic freedom]
Take two generators with \([T^1,T^2]=\ii T^3\).  If \(A_\mu=A_\mu^1T^1\) and \(A_\nu=A_\nu^2T^2\), then the commutator part of \(F_{\mu\nu}\) contains
\[
  \ii g[A_\mu,A_\nu]=\ii gA_\mu^1A_\nu^2[T^1,T^2]=-gA_\mu^1A_\nu^2T^3.
\]
Two gauge components have produced a third.  That is the algebraic origin of gauge-boson self-coupling.  In Faddeev-Popov Ghosts and BRST Symmetry, section 8, the useful habit is to name the object, the operation, and the assumption before using the compact notation.
\end{example}



\begin{exercise}
Set \(f^{abc}=0\) in the non-Abelian field strength.  What remains?  Use the notation of Faddeev-Popov Ghosts and BRST Symmetry: Review problems and extensions: reading asymptotic freedom.
\begin{answer}
The Abelian field strength \(\p_\mu A_\nu^a-\p_\nu A_\mu^a\) for each component.
\end{answer}
\end{exercise}

\begin{exercise}
Why do ghosts interact in Yang-Mills theory?  Use the notation of Faddeev-Popov Ghosts and BRST Symmetry: Review problems and extensions: reading asymptotic freedom.
\begin{answer}
The Faddeev-Popov operator contains the covariant derivative, which contains \(A_\mu\).
\end{answer}
\end{exercise}



\section{Cumulative worked derivation: from deriving the non-Abelian gauge transformation to checking BRST nilpotency}

This final worked section braids together deriving the non-Abelian gauge transformation, varying the Yang-Mills action, and checking BRST nilpotency for Faddeev-Popov Ghosts and BRST Symmetry.  The vocabulary of the chapter was ghosts, gauge volume, BRST, physical states; here those words are used in one continuous calculation rather than in isolated fragments.

Repeat the Abelian construction with a matrix \(U(x)\).  Because matrices do not commute, the commutator of covariant derivatives contains \([A_\mu,A_\nu]\).  That term makes \(F_{\mu\nu}\) nonlinear in \(A_\mu\), and the Yang-Mills action contains cubic and quartic gauge-field interactions.  Gauge fixing the path integral produces a determinant; representing that determinant by anticommuting fields gives Faddeev-Popov ghosts.  In Faddeev-Popov Ghosts and BRST Symmetry, cumulative derivation, the useful habit is to name the object, the operation, and the assumption before using the compact notation.

For reference, the compact formulas are

\[
  D_\mu=\p_\mu+\ii gA_\mu^aT^a
\]
\[
  F_{\mu\nu}^a=\p_\mu A_\nu^a-\p_\nu A_\mu^a-gf^{abc}A_\mu^bA_\nu^c
\]
\[
  \Lag_{\mathrm{YM}}=-\frac14F_{\mu\nu}^aF^{a\mu\nu}
\]

\begin{exercise}
Reproduce the cumulative derivation for Faddeev-Popov Ghosts and BRST Symmetry without looking at the prose.  Mark the step where a finite calculation becomes a continuum, operator, or field-theoretic statement.
\begin{answer}
The reconstruction should name the starting object, carry out the differentiating, varying, transforming, or commuting step explicitly, and state the normalization or boundary condition that makes the final formula meaningful.
\end{answer}
\end{exercise}



\section{Chapter synthesis}

This chapter has treated faddeev-popov ghosts and brst symmetry as a chain of calculations.  The safest summary is not a list of final formulas, but a map from assumptions to equations: finite model, limiting step, boundary condition, normalization, and field-theory use.  If one of those links is missing, the formula should be rederived before it is used in a later chapter.

\begin{exercise}
Write a one-page reconstruction of faddeev-popov ghosts and brst symmetry.  Include one equation from the finite model, one equation from the continuum or operator form, and one sentence explaining how the two are connected.
\begin{answer}
A strong reconstruction names the basic objects, identifies the central derivative, variation, commutator, or symmetry operation, and states the assumption that allows the finite calculation to become the field-theory formula.
\end{answer}
\end{exercise}
